\begin{textsample}{NEG Dim 9 – global\_south\_2024\_02 – Score -1.00 – global\_south\_2024\_02/106253544.txt}  \label{ex:f9_neg_360}
Mediator Qatar said Tuesday it was"hopeful"a new truce between Israel and Hamas could be secured within days, after US President Joe Biden said a ceasefire could start next week and last through Ramadan.
As a dire humanitarian crisis unfolds in the war-battered Gaza Strip, the United Nations humanitarian agency OCHA and the United States separately called to ensure aid reaches the many Palestinians in need.
In the protracted bid to broker a ceasefire nearly five months into the devastating war, Egyptian, Qatari and US mediators have been putting proposals to the parties.
Negotiators are seeking a six-week halt to the fighting and the release of Israeli hostages held in Gaza since Hamas ' s October 7 attack on southern Israel sparked the war.
Several hundred Palestinian detainees held by Israel may also be released under the deal, media reports suggest."My hope is by next Monday we ' ll have a ceasefire"but"we ' re not done yet", Biden said.
Qatar ' s emir, Sheikh Tamim bin Hamad Al-Thani -- whose country hosts Hamas @ @ @ @ @ @ @ @ @ @ November -- met in Paris with French President Emmanuel Macron.
Qatari foreign ministry spokesman Majed al-Ansari said Doha was"hopeful, not \textbf{necessarily} optimistic, that we can announce something"before Thursday."We are going to push for a pause before the beginning of Ramadan", the Muslim fasting month which starts on March 10 or 11, depending on the lunar calendar, Ansari said."We are all aiming towards that target, but the situation is still fluid on the ground."There has been huge international pressure for Israel to hold off on sending troops into Rafah, where nearly 1.5 million Palestinian civilians have sought refuge.
Israeli Prime Minister Benjamin Netanyahu has stressed that any truce would delay, not prevent, a ground invasion of Rafah in Gaza ' s far south, which he said was necessary to achieve"total victory"over Hamas. ' Nothing but dust ' Israel ' s military campaign has killed at least 29,878 people in Gaza, mostly women and children @ @ @ @ @ @ @ @ @ @ Cairo ' s Foreign Minister Sameh Shoukry warned that any assault on Rafah, a key entry point for aid on Gaza ' s border with Egypt,"would have catastrophic repercussions"across the region.
Ahead of the threatened ground incursion, the area has been hit repeatedly by Israeli air strikes.
Abu Khaled Zatmeh, whose nephews were killed in bombardment in Rafah, said the family was"preparing food and getting ready to eat when the strike occurred, and three floors collapsed suddenly"."People began pulling out the martyrs, all of whom were my nephews,"he told AFP.
He said there were"no more"basic supplies in the besieged territory, adding that"even if they allow people to return to the north, there are no houses left -- nothing but dust".
Meanwhile in northern Gaza, desperate Palestinians have scavenged for food, with many people eating animal fodder and even leaves."I have not eaten for two days,"said @ @ @ @ @ @ @ @ @ @ the north, where children roamed with empty pots in search of food."There is nothing to eat or drink."Aid ' under fire ' Most aid trucks have been halted, but foreign armies have air dropped supplies including on Tuesday over Rafah and Gaza ' s main southern city Khan Yunis.
OCHA spokesman Jens Laerke said aid convoys headed for northern Gaza"have come under fire and are systematically denied access to people in need".
The main UN aid agency for Palestinians, UNRWA, said humanitarian relief entering Gaza has halved in February from the previous month, with the latest aid convoy allowed into the north on January 23.
The US Agency for International Development announced another \$53 million in humanitarian assistance for Palestinians, with its head Samantha Power saying aid"has to reach people in need".
Humanitarian workers on the ground"have to be protected", she said on a visit to Jordan.
Along Israel ' s northern border with Lebanon, which @ @ @ @ @ @ @ @ @ @ Hamas ally Hezbollah, UN peacekeepers reported"an expansion and intensification of strikes".
As fierce hostilities continued Tuesday a day after a rare Israeli strike in eastern Lebanon, far from the border, UNIFIL urged deescalation to"leave space to a political and diplomatic solution".
Violence has also surged in the occupied West Bank, where Israeli troops killed three Palestinians in an overnight raid on the Faraa refugee camp.
ABOUT US ARY NEWS brings you 24/7 Live Streaming, Headlines, Bulletins, Talk Shows, Infotainment, and much more.
Watch minute-by-minute updates of current affairs and happenings from Pakistan and all around the world!

% matched lemmas: necessarily
\end{textsample}
